\section{Introduction}
\label{sec:introduction}

Agent-based models (ABMs) represent a system as a population of autonomous agents that follow
local rules and interact through a shared environment \citep{macal2010tutorial}. Their appeal
is that aggregate patterns, such as an epidemic curve, a segregated city
\citep{schelling1971dynamic} or an unequal wealth distribution, emerge from rules a domain
expert can state and criticise. Their weakness is the cost of answering questions with them.
A single run tells little: a credible study runs many replicates per parameter combination,
explores the parameter space, fits the parameters to data, asks which of them matter, and
reports all of it in a form others can rerun. Each of these steps multiplies the number of
runs, and each is usually assembled by hand around the simulation library.

Python has become a common language for this work, and Mesa \citep{kazil2020mesa,
terhoeven2025mesa3} its best-known ABM framework. Python makes models easy to write and to
connect to the scientific stack \citep{harris2020numpy, virtanen2020scipy}, but per-agent
Python code is slow, and the surrounding tasks (sweeps, calibration, sensitivity analysis,
provenance) are left to the modeller or to separate tools \citep{herman2017salib,
klinger2018pyabc, kwakkel2017ema}. Faster platforms exist in Java and Julia
\citep{luke2005mason, north2013repast, datseris2024agents} and on GPUs
\citep{richmond2023flamegpu}; in Python, mesa-frames trades the per-agent programming model for
vectorised speed \citep{mesaframes}, and Repast4Py distributes a run over MPI processes
\citep{collier2020repast4py}.

\modelly{} is a new framework, written from scratch, that removes this cost without asking the
modeller to leave Python or to give up agents. Its programming model is compatible with
Mesa's, so a Mesa user is at home from the first line and a codemod migrates an existing Mesa
model's framework calls. A researcher writes the model once, as plain Python with agents as
objects, and gets:

\begin{itemize}[leftmargin=*,itemsep=2pt]
  \item \textbf{Speed without leaving Python.} Agent attributes can live in model-owned
  \code{numpy} arrays (\emph{columns}), identical agents can be counted rather than
  instantiated (\emph{cohorts}), and the framework's inner loops run in compiled C, each with a
  pure-Python twin that gives the same results (\cref{sec:design}). On eight workloads the
  column form runs 4.6 to 42 times faster than Mesa~2.3 and 2.9 to 33 times faster than
  Mesa~3.5, the cohort form 251 and 135 times, and an unchanged Mesa model runs up to 4.8 times
  faster; a study of 100 seeds per variant shows that the faster forms compute the same model
  (\cref{sec:performance}).
  \item \textbf{The whole study in one package.} Parameter sweeps, approximate Bayesian
  calibration, Morris, Sobol and PAWN sensitivity analysis, emulators and history matching,
  provenance manifests, verifiable archive bundles and ODD descriptions \citep{grimm2006odd,
  grimm2020odd} are part of the framework, each a few lines of code, and share one seeding
  discipline, so a study's results do not depend on the executor, the number of workers or the
  order in which runs finish (\cref{sec:workflow}). Each method is checked against problems
  whose answers are known (\cref{sec:reproducibility}).
  \item \textbf{Results others can rerun.} Seeded generators reproduce the standard library's
  random stream; a run's manifest records everything needed to repeat it, and
  \cmd{modelly verify} repeats it and names the first output that differs; a determinism suite
  checks that the bundled examples give the same digests compiled or in pure Python and
  straight or resumed from a checkpoint (\cref{sec:reproducibility}).
  \item \textbf{From one laptop to large and data-driven studies.} The same model can be split
  across processes with byte-identical results, run its array phases on a GPU (35 times faster
  at ten million agents), be differentiated for gradient calibration, assimilate live
  observations, couple agents to diffusing fields, be emulated when runs are expensive, start
  from a synthetic population, and be shared as a page that runs in any browser
  (\cref{sec:scale}).
\end{itemize}

The rest of the paper situates \modelly{} among existing tools (\cref{sec:related}),
describes its design (\cref{sec:design}) and research workflow (\cref{sec:workflow}),
evaluates its speed and its agreement with Mesa (\cref{sec:performance}) and the correctness
of its methods (\cref{sec:reproducibility}), measures how it scales out and connects to data
(\cref{sec:scale}), sets out the studies it serves (\cref{sec:usecases}), and ends with the threats to the validity of these measurements and the
software's limitations (\cref{sec:threats}).
